% Threat model and LINDDUN privacy threats for FitBuddy (verified against source, version 3.2.82).
% Required in the preamble:
%   \usepackage{booktabs}
%   \usepackage{tabularx}
%   \newcolumntype{L}{>{\raggedright\arraybackslash}X}

\section{Threat Model}
\label{sec:threat-model}

The threat model is derived from the architecture in Figure~\ref{fig:architecture} and was
verified against the FitBuddy source code (version 3.2.82). We first describe the data flows
that leave the application sandbox, then the information at risk, our trust assumptions, the
trust boundaries, the potential adversaries, and finally the data protection roles. The privacy
threats that follow from this model are analysed and prioritised in
Section~\ref{sec:privacy-threats}.

\subsection{Data Flows}
FitBuddy is designed as a local-first application: there is no FitBuddy account and no central
application server. Food, workout and body-measurement records are stored in an on-device Room
(SQLite) database; settings, including the user's name and AI API keys, are stored in a
Preferences DataStore. Meal photos are downscaled and re-encoded before upload, which strips EXIF
metadata such as GPS coordinates, and are not stored in the database. Barcodes are decoded on the
device. Nevertheless, personal data leaves the sandbox through the eight flows in
Table~\ref{tab:flows}.

\begin{table}[htbp]
\centering
\footnotesize
\caption{Data flows leaving the FitBuddy application sandbox.}
\label{tab:flows}
\begin{tabularx}{\textwidth}{@{}l L L L@{}}
\toprule
\textbf{ID} & \textbf{Data} & \textbf{Destination} & \textbf{Trigger / default} \\
\midrule
F1 & Meal photo or text, plus a \emph{user-state context}: age, sex, height, weight, goal,
     activity level, calorie targets, today's intake and expenditure, and the latest body
     measurement (weight, body fat, BMR, muscle mass)
   & User-selected AI provider: OpenRouter, Google Gemini, Ollama (local or cloud) or an
     OpenAI-compatible host
   & Every AI-assisted food or workout log \\
F2 & First name, profile, nutrition targets and their history, 30 days of daily nutrition,
     exercise and body data, monthly averages of all older history, and the chat conversation
   & Same AI provider
   & Progress insights and AI coach chat \\
F3 & Entire app-private storage: database and DataStore, including name, support ID and API
     keys
   & Google account (Android Auto Backup)
   & Automatic, by the OS, when device backup is on \\
F4 & AES-256-GCM-encrypted export of the whole database and settings (including API keys);
     in plaintext: support ID, chunk IDs, user-defined device name, MAC address or
     \texttt{ANDROID\_ID}
   & Developer-operated proxy on Vercel $\rightarrow$ MongoDB Atlas
   & Opt-in (off by default); once enabled, uploads automatically on app start, at most every
     12\,h \\
F5 & Daily and update \emph{heartbeats}: first and last name, signup date, record count, AI
     provider, device manufacturer and model; crash reports with the support ID as user ID
   & Sentry, directly or via the same Vercel proxy when Sentry is unreachable
   & Heartbeats: always, in all builds. Crash reports: on by default in GitHub builds
     (opt-out on first launch), off in F-Droid builds \\
F6 & All body measurements (read) and new measurements (write)
   & FreeScale companion app via an exported \texttt{ContentProvider}
   & On request of the companion app \\
F7 & Scanned barcodes, IP address; OAuth authorisation code
   & Open Food Facts; jsDelivr and GitHub (exercise catalog, emoji, donor list, update check);
     F-Droid; OpenRouter OAuth via the browser and a loopback server on \texttt{127.0.0.1}
   & Barcode scan, catalog refresh, app start, OpenRouter sign-in \\
F8 & Backup file (plaintext JSON unless a password is set) or diagnostic log
   & Any app chosen by the user in the Android share sheet
   & Manual export \\
\bottomrule
\end{tabularx}
\end{table}

\subsection{Information at Risk}
Table~\ref{tab:assets} lists the assets processed by FitBuddy. The most important property is
\emph{aggregation}: a single record is of limited value, but together the records form a
longitudinal health profile that can be linked to a named person through the identifiers in A4.

\begin{table}[htbp]
\centering
\footnotesize
\caption{Information at risk.}
\label{tab:assets}
\begin{tabularx}{\textwidth}{@{}l l L l L@{}}
\toprule
\textbf{ID} & \textbf{Asset} & \textbf{Examples} & \textbf{Flows} & \textbf{Sensitivity} \\
\midrule
A1 & Health \& body data
   & Age, sex, height, weight, goals; smart-scale readings such as body fat, BMR, muscle mass,
     visceral fat and metabolic age
   & F1--F4, F6, F8 & Special category (GDPR Art.~9) \\
A2 & Dietary record
   & Food logs, ingredients, saved meals, meal photos, scanned products
   & F1--F4, F7, F8
   & High: can reveal illness, eating disorders, pregnancy or religion (e.g.\ halal, fasting);
     photos may show surroundings or other people \\
A3 & Activity data & Workouts, sessions, calories burned & F1--F4, F8 & Health data (Art.~9) \\
A4 & Identity \& identifiers
   & First and last name, support ID (persistent UUID), user-defined device name,
     \texttt{ANDROID\_ID} or MAC, signup date
   & F2--F5 & High: links all other assets to a person, across services \\
A5 & Credentials
   & AI provider API keys, OpenRouter OAuth key, cloud-backup password (encrypted with the
     Android Keystore)
   & F1, F3, F4, F8
   & High: abuse of paid keys; access to provider-side request history \\
A6 & Inferred \& conversational data
   & Region, AI coach conversations, progress insights
   & F1--F4
   & High: region can indicate cultural or ethnic background (Art.~9); enables profiling \\
\bottomrule
\end{tabularx}
\end{table}

\subsection{Trust Assumptions}
\begin{description}
  \item[TA1 -- Android platform.] The OS, the application sandbox and the Android Keystore work as
    specified and the device is not rooted. Other apps can therefore reach FitBuddy's data only
    through components that the app deliberately exports.
  \item[TA2 -- Transport security.] TLS to HTTPS endpoints is sound and validated against the
    system CA store; no certificate pinning is used. This assumption does \emph{not} hold for
    user-configured local endpoints, because cleartext HTTP is permitted application-wide
    (\texttt{usesCleartextTraffic="true"}) to support self-hosted models on the LAN.
  \item[TA3 -- The user.] The user controls the device and their Google account and chooses an
    AI provider deliberately. We do \emph{not} assume that the user knows which profile data is
    attached to each AI request, or that heartbeats are sent.
  \item[TA4 -- AI providers are semi-trusted.] Providers perform the requested inference correctly
    (honest-but-curious), but retention, logging, human review and training use are governed by
    their own terms, which often permit training on free tiers. OpenRouter is a broker that
    forwards requests to further model hosts.
  \item[TA5 -- Developer infrastructure is semi-trusted.] The Sentry project, the Vercel proxy and
    the MongoDB Atlas cluster are operated by the developer. We do not assume malicious intent,
    but the developer can read anything sent in plaintext, and the infrastructure can be breached.
    Its access control relies on one API key that is shared by all installations and embedded,
    obfuscated, in every APK, so it must be considered public.
  \item[TA6 -- Google backup.] The confidentiality of F3 depends on the security of the user's
    Google account; backups are only end-to-end encrypted when the device has a screen lock.
  \item[TA7 -- Out of scope.] Compromise of the build pipeline, third-party libraries or signing
    keys, and attacks on the language models themselves (e.g.\ prompt injection that corrupts
    estimates).
\end{description}

\subsection{Trust Boundaries}
The architecture contains seven trust boundaries. Several of them are crossed
\emph{without user interaction} for each transfer (F3, F4 once enabled, F5).

\begin{description}
  \item[TB1 -- App sandbox $\leftrightarrow$ other apps on the device (F6, F8).] The measurement
    bridge is an exported \texttt{ContentProvider} without an Android permission; it authorises
    callers only by comparing their package name to an allow-list. Manual exports are handed to
    whichever app the user picks.
  \item[TB2 -- Device $\leftrightarrow$ Google cloud (F3).] Crossed by the OS without application
    control. The backup rules deliberately include all databases and preferences, including API
    keys, so that a reinstall restores them.
  \item[TB3 -- Device $\leftrightarrow$ AI provider (F1, F2).] The main boundary for health data:
    each request carries profile and health context, not only the photo or text that needs
    estimating. For LAN providers, the local network also lies on this boundary.
  \item[TB4 -- Device $\leftrightarrow$ developer infrastructure (F4, F5).] When the Sentry host is
    unreachable (e.g.\ because of a DNS-based blocker), telemetry is rerouted through the Vercel
    proxy, so the user's own network controls do not enforce this boundary.
  \item[TB5 -- Device $\leftrightarrow$ public services (F7).] Only query metadata crosses, but
    barcodes reveal which products a user consumes, and every request reveals the user's IP
    address.
  \item[TB6 -- Device $\leftrightarrow$ person holding it.] Data at rest has no application-level
    encryption (no SQLCipher, no encrypted preferences) and relies only on the OS's file-based
    encryption while the device is locked.
  \item[TB7 -- Support-ID holder $\leftrightarrow$ cloud backup (F4).] Backups are stored and
    retrieved under the support ID (\texttt{/api/backup/\{supportId\}}). Since the proxy key is
    public (TA5), \emph{knowing a support ID is sufficient to download that user's backup}. The
    app tells users to share this ID when reporting a bug and to keep it safe for cloud restore,
    so it acts both as a public identifier and as a credential.
\end{description}

\subsection{Potential Adversaries}
Table~\ref{tab:adversaries} relates each adversary to the boundary it acts on and to the assets
it can reach.

\begin{table}[htbp]
\centering
\footnotesize
\caption{Adversaries, their capabilities, and the trust boundaries they attack.}
\label{tab:adversaries}
\begin{tabularx}{\textwidth}{@{}l L L l l@{}}
\toprule
\textbf{ID} & \textbf{Adversary} & \textbf{Capability} & \textbf{TB} & \textbf{Assets} \\
\midrule
ADV1 & Network attacker (public Wi-Fi, malicious LAN)
     & HTTPS: metadata only (destination, timing, upload size, which reveal meal times).
       Cleartext LAN endpoints: read and modify photos, health context and Bearer API keys
     & TB3, TB5 & A1--A3, A5 \\
ADV2 & Malicious app on the device
     & If installed under the package name \texttt{com.anant.freescale} while the genuine
       companion app is absent, it passes the allow-list check, exports all measurements and can
       inject falsified ones
     & TB1 & A1 \\
ADV3 & Semi-trusted AI provider
     & Receives every photo plus health context and, for the coach, months of history and the
       first name; can retain, train on and link data to the account behind the API key
     & TB3 & A1--A4, A6 \\
ADV4 & Curious service operator (developer)
     & Receives name, signup date and device data in daily heartbeats; holds all cloud backups
       and can decrypt those made without a custom password, because the key is derived from
       the support ID that is sent with them
     & TB4, TB7 & A1--A6 \\
ADV5 & Attacker breaching developer infrastructure
     & Same as ADV4, but for all users at once (Sentry, Vercel, MongoDB Atlas)
     & TB4 & A1--A6 \\
ADV6 & Support-ID holder (anyone given the ID for a bug report, or who sees it)
     & Extracts the proxy key from the APK, downloads the backup by support ID, and decrypts it
       if no custom password was set
     & TB7 & A1--A6 \\
ADV7 & Person with physical access (partner, thief, forensic examiner)
     & Unlocked device: plaintext database, API keys and unencrypted exports
     & TB6, TB1 & A1--A6 \\
ADV8 & Google account attacker
     & Restores the Auto Backup on another device, obtaining health data and API keys
     & TB2 & A1--A6 \\
ADV9 & Public service observer (Open Food Facts, jsDelivr, GitHub)
     & Links the user's IP address to scanned products and app usage
     & TB5 & A2 \\
\bottomrule
\end{tabularx}
\end{table}

Two observations shape the analysis that follows. First, the most capable adversaries are not
external attackers but parties that receive data \emph{by design} (ADV3, ADV4) or through an
identifier that users are invited to share (ADV6). Their risk follows from design decisions
rather than from implementation bugs. Second, the classic network eavesdropper is largely
mitigated by TLS; the realistic network risk is confined to user-configured cleartext endpoints.

\subsection{Data Protection Roles}
Under the GDPR, the developer acts as \emph{controller} for F4 and F5, since they determine the
purpose and means of the telemetry and backup infrastructure; Sentry, Vercel and MongoDB act as
their processors. For F1 and F2 the responsibility is split: users bring their own API keys and
therefore contract with the AI provider directly, while the application decides which personal
data is attached to each request. Google processes F3 under the user's own Google account
agreement. Because A1, A3 and possibly A6 are special-category data, processing in F1, F2, F4 and
F5 requires explicit consent under Art.~9(2)(a) GDPR, and the developer is responsible for the
accuracy of the information given to users (Art.~13).

\section{Privacy Threats}
\label{sec:privacy-threats}

We elicited privacy threats by applying the LINDDUN categories (Linkability, Identifiability,
Non-repudiation, Detectability, Disclosure of information, Unawareness and Non-compliance) to each
data flow in Table~\ref{tab:flows}.

\paragraph{Prioritisation criteria.} Each threat receives two scores from 1 (low) to 3 (high):
\begin{itemize}
  \item \textbf{Impact:} the sensitivity of the data exposed (Art.~9 data scores highest), the
        number of users affected, and whether the exposure contradicts what users are told.
  \item \textbf{Likelihood:} whether the exposure happens by default or needs an opt-in or
        unusual configuration, and how much effort the adversary needs.
\end{itemize}
The priority is the product of both scores. Ties are broken by the number of affected users.

\begin{table}[htbp]
\centering
\footnotesize
\caption{Privacy threats, ordered by priority (I = impact, L = likelihood, P = I $\times$ L).
LINDDUN: L = linkability, I = identifiability, D = disclosure of information,
U = unawareness, Nc = non-compliance.}
\label{tab:threats}
\begin{tabularx}{\textwidth}{@{}l L l l c c c@{}}
\toprule
\textbf{ID} & \textbf{Threat} & \textbf{LINDDUN} & \textbf{Flow / Adv.} & \textbf{I} & \textbf{L} & \textbf{P} \\
\midrule
T1 & Daily heartbeats send the user's full name, signup date and usage counts to Sentry in all
     builds, even after opting out of crash reporting. The onboarding screen states that the
     user's name is never sent.
   & U, Nc, I & F5 / ADV4--5 & 2 & 3 & 6 \\
T2 & Every AI request carries health context that is not needed for the task; the coach
     additionally sends the first name and months of history to a provider that may retain or
     train on it
   & D, Nc & F1, F2 / ADV3 & 3 & 2 & 6 \\
T3 & The default backup key is derived from the support ID, which is stored next to the
     ciphertext and shared in bug reports. With the public proxy key, anyone holding the ID can
     download and decrypt the full history and API keys
   & D & F4 / ADV4--6 & 3 & 2 & 6 \\
T4 & Persistent identifiers (support ID, \texttt{ANDROID\_ID}, user-defined device name) link
     telemetry, backups and the public donor list to one person
   & L, I & F4, F5 / ADV4--5 & 2 & 2 & 4 \\
T5 & Telemetry is rerouted through the developer's proxy when the user blocks Sentry, overriding
     the user's network-level choice
   & U, Nc & F5 / ADV4 & 2 & 2 & 4 \\
T6 & Google Auto Backup copies all health data and API keys to the cloud without informing the
     user
   & D, U & F3 / ADV8 & 3 & 1 & 3 \\
T7 & No application-level encryption at rest; manual exports are plaintext JSON by default
   & D & F8 / ADV7 & 3 & 1 & 3 \\
T8 & Cleartext HTTP is allowed app-wide, exposing photos, health context and API keys to the
     local network when a LAN provider is used
   & D & F1 / ADV1 & 3 & 1 & 3 \\
T9 & The measurement bridge verifies callers by package name only, not by signing certificate
   & D & F6 / ADV2 & 2 & 1 & 2 \\
T10 & Barcode lookups link the user's IP address to the products they eat
   & L & F7 / ADV9 & 1 & 2 & 2 \\
\bottomrule
\end{tabularx}
\end{table}

The three highest-priority threats share a pattern: they result from design choices rather than
coding errors, and each conflicts with a statement or expectation given to the user. They are
discussed in detail in the following sections.
